\begin{afterword}
\summaryhead

Surprised that commenters defended the Allais choices, the author appeals to two lessons from the study of heuristics and biases: subjects defend incoherent preferences, and people put a large extra value on certainty.

The first is illustrated by a recalled preference reversal, in which people choose one bet but price the other higher and so can be pumped for money, and by a transcript of a Las Vegas subject who keeps such a pattern. The second is based on a recalled experiment. The trouble with valuing certainty, the post argues, is that one time's certainty is another time's probability: a chain of events can remove probability step by step until a preference flips. It restates the Allais money pump and gives a framing example.

The post then tells the reader to multiply utilities by probabilities, with a rule for the Allais problem under logarithmic utility. Against ``the utility of certainty'' it says that utilities attach to outcomes, not probabilities. Paying for a feeling of certainty is human, but when lives are at stake, preferences that reverse show that one is not steering toward any goal.

\responsehead

The post's normative point is the one Kahneman and Tversky themselves made: the departures from expected utility that their theory describes lead to inconsistencies. Its framing example works as stated, and the design it describes was tested in their 1979 paper, with the result the post predicts.

The evidence for the two lessons is recalled rather than cited, and the post says so (``IIRC,'' ``I believe''). The two bets for the preference reversal could not be traced. Preference reversal itself, in which people choose one bet but put a higher price on the other, is well established; it was found by Lichtenstein and Slovic, and Kahneman and Tversky note that their theory does not treat bidding. The recalled certainty result, that a drop from 100 to 99 percent outweighed one from 80 to 20 percent, is far stronger than the nearest published finding, which compares a drop from 1.0 to .25 with one from .8 to .2. The claim that the certainty effect is ``even stronger over human lives'' has no source; the best-known study with lives found that a sure option was preferred when outcomes were described as lives saved, and a gamble when they were described as deaths.

On persistence, the evidence points both ways. Studies of the Allais pattern have found violations that survive arguments for the axiom and also a tendency to correct them. The post reads persistence as a failure to see an error; others read the same finding as support for Allais.

The answer to ``the utility of certainty'' is right against the commenter's formula, which puts a probability inside a utility function. The published versions of the objection do something else: they redescribe the outcomes, counting the disappointment of losing a sure thing, or the way a sum was obtained, as part of the outcome. The post answers the version about feelings; the other it does not discuss.

The worked rule for logarithmic utility contains a slip. By the post's own method, 1A is better than 1B only for assets under about \$550, not for anyone whose assets \$24,000 would double. The slip leaves the post's point intact.

In short: a restatement of the case against the Allais choices, whose recalled experiments are marked as recalled, and whose answer to ``the utility of certainty'' answers a commenter's formula but not the published forms of that objection.
\end{afterword}
